\chapter{Gaya dan Asas Kelembaman}\label{ch:g10:inertia}

Berhentilah mengayuh dan sepeda itu meluncur sampai berhenti;
berhentilah menarik dan kereta luncur itu terdiam. Pengalaman berbisik
bahwa gerak harus disuapi usaha --- dan selama dua ribu tahun fisika
menyetujuinya. Bab ini membangun \omterm{def:g10:inertia:force}{gaya}, sebuah anak panah yang diukur
dalam \omterm{def:g10:inertia:force}{newton}, lalu menyatakan hukum besar pertama mekanika: apabila
benar-benar dibiarkan sendirian, benda yang bergerak mempertahankan
geraknya, lurus dan tetap, selamanya.

\section{Memodelkan sebuah pengaruh sebagai gaya}

\begin{definition}[Gaya]\label{def:g10:inertia:force}
Sebuah \emph{gaya}\index{gaya} memodelkan pengaruh mekanis (dorongan,
tarikan, tarikan gravitasi) yang dikerjakan pada sebuah benda. Gaya
dicirikan oleh:
\begin{itemize}
  \item \emph{titik tangkapnya}\index{titik tangkap};
  \item arahnya --- garis tempat gaya itu bekerja beserta arah pada
        garis tersebut;
  \item besarnya, yang diukur dalam \emph{newton}\index{newton}
        (lambang \unit{N}).
\end{itemize}
Sebuah gaya digambarkan sebagai anak panah $\vect F$ (sebuah vektor,
seperti pada jilid matematika) dari titik tangkapnya, dengan panjang
yang sebanding dengan besarnya pada skala yang disebutkan.
\end{definition}

\begin{example}[Berat sebuah tas sekolah]\label{ex:g10:inertia:schoolbag}
Sebuah tas sekolah bermassa $m = \qty{4.0}{kg}$ ditarik ke bawah oleh
Bumi dengan \omterm{def:g10:universal-gravitation:weight}{beratnya} (\cref{ch:g10:universal-gravitation}):
$P = mg = 4.0 \times 9.81 \approx \qty{39}{N}$ --- satu \omterm{def:g10:inertia:force}{newton} kira-kira
\omterm{def:g10:universal-gravitation:weight}{berat} sebuah apel. Pada skala $\qty{1}{cm} \leftrightarrow
\qty{10}{N}$: sebuah anak panah tegak sepanjang \qty{3.9}{cm}, bertitik
tangkap di titik \omterm{def:g10:universal-gravitation:weight}{berat} tasnya, menunjuk ke bawah.
\end{example}

\section{Daftar gaya}

\begin{definition}[Gaya sentuh dan gaya jarak jauh]\label{def:g10:inertia:contact}
\emph{Gaya sentuh}\index{gaya sentuh} dikerjakan pada titik tempat dua
benda bersentuhan. \emph{Gaya jarak jauh}\index{gaya jarak jauh} tidak
memerlukan sentuhan: \omterm{def:g10:universal-gravitation:weight}{berat} --- tarikan gravitasi pada
\cref{ch:g10:universal-gravitation} --- adalah satu-satunya \omterm{def:g10:inertia:force}{gaya} jarak
jauh yang dipakai pada bab ini.
\end{definition}

\begin{definition}[Gaya sentuh yang lazim]\label{def:g10:inertia:inventory}
Empat \omterm{def:g10:inertia:contact}{gaya sentuh} sudah mencakup hampir semua keadaan sehari-hari:
\begin{itemize}
  \item \emph{gaya normal}\index{gaya normal} $\vect R$ dari sebuah
        tumpuan, tegak lurus permukaannya, mendorong menjauhinya;
  \item \emph{tegangan tali}\index{tegangan tali} $\vect T$ pada kawat,
        tali atau rantai, terarah sepanjang tali itu, menarik ke
        arahnya;
  \item \emph{gaya gesek}\index{gaya gesek} $\vect f$ dari sebuah
        permukaan, udara atau air, yang melawan peluncuran atau
        geraknya;
  \item \emph{gaya dorong}\index{gaya dorong} mesin, semburan atau
        tangan, yang mendorong searah usahanya.
\end{itemize}
\end{definition}

\begin{omfigure}
\begin{tikzpicture}[scale=0.9, font=\small]
  \draw[fill=gray!15] (-2.4,0) rectangle (2.4,-0.25);
  \draw (-2.0,-0.25) -- (-2.0,-1.5);
  \draw (2.0,-0.25) -- (2.0,-1.5);
  \draw[fill=omExo!25] (-0.9,0) rectangle (0.9,0.5);
  \coordinate (G) at (0,0.25);
  \fill (G) circle (1.5pt);
  \draw[-Stealth, omThm, very thick] (G) -- (0,-1.3)
    node[right] {$\vect P$};
  \draw[-Stealth, omDef, very thick] (G) -- (0,1.8)
    node[right] {$\vect R$};
\end{tikzpicture}

{\small Buku di atas meja: \omterm{def:g10:universal-gravitation:weight}{berat} $\vect P$ dan \omterm{def:g10:inertia:inventory}{gaya normal} $\vect R$ ---
garis yang sama, panjang yang sama, arah yang berlawanan.}
\end{omfigure}

\begin{example}[Daftar gaya pada kereta luncur yang ditarik]\label{ex:g10:inertia:sled}
Sebuah kereta luncur ditarik melintasi salju: \omterm{def:g10:universal-gravitation:weight}{berat} $\vect P$ (jarak
jauh), \omterm{def:g10:inertia:inventory}{gaya normal} $\vect R$ dari saljunya, tegangan tali $\vect T$ pada
talinya, \omterm{def:g10:inertia:inventory}{gaya gesek} $\vect f$ yang melawan peluncurannya. Tidak ada lagi
yang menyentuh kereta itu: daftarnya lengkap --- \emph{\omterm{def:g10:universal-gravitation:weight}{berat} selalu ada,
lalu satu \omterm{def:g10:inertia:contact}{gaya sentuh} untuk tiap benda yang bersentuhan}.
\end{example}

\section{Asas kelembaman}

\begin{definition}[Gerak lurus beraturan]\label{def:g10:inertia:uniform}
Sebuah benda berada dalam \emph{gerak lurus beraturan}\index{gerak lurus
beraturan} di dalam sebuah kerangka apabila \omterm{def:g10:relative-motion:trajectory}{lintasannya} berupa garis
lurus yang ditempuh dengan kelajuan tetap: vektor kecepatannya $\vect v$
sama pada setiap saat. Keadaan diam adalah kasus khusus
$\vect v = \vect 0$.
\end{definition}

\begin{definition}[Gaya yang saling meniadakan]\label{def:g10:inertia:compensate}
\omterm{def:g10:inertia:force}{Gaya} yang bekerja pada sebuah benda dikatakan \emph{saling
meniadakan}\index{gaya saling meniadakan} apabila pengaruh gabungannya
nol: disusun ujung ke pangkal, anak panahnya kembali ke titik semula.
Dua \omterm{def:g10:inertia:force}{gaya} saling meniadakan tepat apabila keduanya berbagi garis kerja
yang sama dengan besar yang sama dan arah yang berlawanan.
\end{definition}

\begin{theorem}[Asas kelembaman]\label{thm:g10:inertia:principle}
\index{asas kelembaman}\index{kelembaman}
Di dalam \omterm{def:g10:relative-motion:frames}{kerangka tanah}, benda yang tidak dikenai \omterm{def:g10:inertia:force}{gaya}, atau dikenai
\omterm{def:g10:inertia:force}{gaya} yang \omterm{def:g10:inertia:compensate}{saling meniadakan}, tetap diam atau mempertahankan \omterm{def:g10:inertia:uniform}{gerak lurus
beraturannya}; sebaliknya, benda yang diam atau bergerak lurus beraturan
dikenai \omterm{def:g10:inertia:force}{gaya} yang \omterm{def:g10:inertia:compensate}{saling meniadakan} (atau tidak dikenai \omterm{def:g10:inertia:force}{gaya} sama
sekali).
\end{theorem}
\admitted

\begin{remark}
Tidak ada percobaan yang sanggup mengasingkan sebuah benda sepenuhnya,
sehingga asas ini merupakan postulat --- disimak sekilas oleh Galileo,
dinyatakan oleh Newton. Bab berikutnya meningkatkannya menjadi hukum
yang mengaitkan \omterm{def:g10:inertia:force}{gaya} dengan \emph{perubahan} kecepatan; jilid Tahun ke-1
bertanya di dalam kerangka yang mana asas itu sungguh berlaku.
\end{remark}

\begin{proposition}[Membaca gerak dari gaya, dan sebaliknya]\label{prop:g10:inertia:converse}
Di dalam \omterm{def:g10:relative-motion:frames}{kerangka tanah}:
\begin{enumerate}
  \item jika kecepatan sebuah benda berubah --- baik besarnya
        \emph{maupun} arahnya --- maka \omterm{def:g10:inertia:force}{gaya} yang bekerja padanya
        \emph{tidak} \omterm{def:g10:inertia:compensate}{saling meniadakan};
  \item jika sebuah benda diam atau bergerak lurus beraturan, \omterm{def:g10:inertia:force}{gaya} pada
        daftarnya harus \omterm{def:g10:inertia:compensate}{saling meniadakan} --- dan itulah yang menentukan
        besaran yang belum diketahui.
\end{enumerate}
\end{proposition}

\begin{proof}
Pernyataan pertama adalah kontraposisi
\cref{thm:g10:inertia:principle}, dan yang kedua adalah bagian
sebaliknya yang diterapkan pada daftar \omterm{def:g10:inertia:force}{gaya} yang sudah lengkap.
\end{proof}

\begin{omfigure}
\begin{tikzpicture}[scale=0.85, font=\small]
  \draw[gray!60] (-0.5,0) -- (10.3,0) node[below left, gray] {es};
  \foreach \x in {0.5,2.5,4.5,6.5,8.5} {
    \fill (\x,0.13) ellipse (0.3 and 0.13);
    \draw[-Stealth, omDef, thick] (\x,0.5) -- (\x+1.3,0.5);
  }
  \node[omDef] at (1.15,0.9) {$\vect v$};
  \node[omDef] at (5.4,1.5) {langkah yang sama dalam selang waktu yang sama};
\end{tikzpicture}

{\small Sebuah keping hoki sesudah dilepaskan: kedudukannya pada selang
waktu yang sama, masing-masing dengan satu anak panah kecepatan yang
serupa --- \omterm{def:g10:inertia:uniform}{gerak lurus beraturan}.}
\end{omfigure}

\begin{example}[Keping hoki]\label{ex:g10:inertia:puck}
Setelah dilepaskan tongkatnya, sebuah keping menyeberangi \qty{30}{m}
es yang licin dalam \qty{2.5}{s}, lurus dengan kelajuan tetap
$v = 30/2.5 = \qty{12}{m/s}$. Daftar \omterm{def:g10:inertia:force}{gayanya}: \omterm{def:g10:universal-gravitation:weight}{berat} dan \omterm{def:g10:inertia:inventory}{gaya normal}
(\omterm{def:g10:inertia:inventory}{gesekannya} dapat diabaikan). Keduanya \omterm{def:g10:inertia:compensate}{saling meniadakan}, dan asas itu
meramalkan persis gerak yang teramati: tidak ada yang mendorong keping
itu ke depan --- dan memang tidak diperlukan.
\end{example}

\begin{remark}[Galileo melawan Aristoteles]
Aristoteles mengajarkan bahwa setiap gerak memerlukan penggerak, dan
kehidupan sehari-hari --- dengan \omterm{def:g10:inertia:inventory}{gesekan} pada setiap sentuhan ---
tampaknya menyetujuinya. Galileo membayangkan \omterm{def:g10:inertia:inventory}{gesekannya} dihilangkan:
permukaan yang makin licin memerlukan dorongan yang makin sedikit untuk
\emph{mempertahankan} gerak. Ketekunan bergerak, bukan keadaan diam,
itulah keadaan alaminya.
\end{remark}

\begin{remark}[Penumpang yang terhuyung]
Ketika sebuah bus mengerem mendadak, penumpang yang berdiri terhuyung
ke depan. Tidak ada \omterm{def:g10:inertia:force}{gaya} yang melemparkan mereka: karena kelembaman,
tubuh mereka mempertahankan kecepatannya sementara busnya kehilangan
kecepatannya sendiri. Sabuk pengamanlah yang memasok \omterm{def:g10:inertia:force}{gaya} ke belakang
yang tidak ada pada daftarnya.
\end{remark}

\begin{omfigure}
\begin{tikzpicture}[scale=0.62, font=\small]
  \draw[gray!60] (-0.3,0) -- (7.3,0);
  \foreach \x in {0.7,2.5,4.3,6.1}
    \fill[omDef] (\x,0.22) ellipse (0.34 and 0.22);
  \draw[-Stealth, omThm, thick] (0.7,0.22) -- (0.7,-1.0)
    node[right] {$\vect P$};
  \draw[-Stealth, omDef, thick] (0.7,0.44) -- (0.7,1.65)
    node[right] {$\vect R$};
  \node at (3.6,-1.7) {gaya saling meniadakan: langkahnya tetap sama};
\end{tikzpicture}
\qquad
\begin{tikzpicture}[scale=0.62, font=\small]
  \draw[gray!60] (-0.3,0) -- (7.3,0);
  \foreach \x in {0.7,2.4,3.8,4.9,5.7,6.2}
    \fill[omProp] (\x,0.22) ellipse (0.34 and 0.22);
  \draw[-Stealth, omThm, thick] (0.7,0.22) -- (0.7,-1.0)
    node[right] {$\vect P$};
  \draw[-Stealth, omDef, thick] (0.7,0.44) -- (0.7,1.65)
    node[right] {$\vect R$};
  \draw[-Stealth, omProp, very thick] (0.7,0.22) -- (-0.3,0.22)
    node[above] {$\vect f$};
  \node at (3.6,-1.7) {gesekan tak terimbangi: batunya melambat};
\end{tikzpicture}

{\small Sebuah batu curling pada selang waktu yang sama. Kiri: $\vect P$
dan $\vect R$ \omterm{def:g10:inertia:compensate}{saling meniadakan} --- langkahnya tetap sama. Kanan: masih
tersisa \omterm{def:g10:inertia:inventory}{gaya gesek} $\vect f$ yang tak terimbangi --- langkahnya menyusut,
kecepatannya berubah.}
\end{omfigure}

\section{Gaya yang setimbang: statika dari sebuah gambar}

\begin{definition}[Diagram gaya]\label{def:g10:inertia:diagram}
\emph{Diagram gaya}\index{diagram gaya} sebuah benda menampilkan benda
itu sendirian, dengan setiap \omterm{def:g10:inertia:force}{gaya} yang bekerja \emph{padanya} digambar
sesuai skala dari \omterm{def:g10:inertia:force}{titik tangkapnya}; \omterm{def:g10:inertia:force}{gaya} yang dikerjakan \emph{oleh}
benda itu tidak pernah muncul.
\end{definition}

\begin{method}[Menganalisis keadaan statis]\label{met:g10:inertia:method}
\begin{enumerate}
  \item Pilihlah benda yang hendak ditelaah; bekerjalah di dalam
        \omterm{def:g10:relative-motion:frames}{kerangka tanah}.
  \item Daftarkan \omterm{def:g10:inertia:force}{gayanya}: \omterm{def:g10:universal-gravitation:weight}{berat} selalu ada, satu \omterm{def:g10:inertia:contact}{gaya sentuh} untuk tiap
        sentuhan.
  \item Gambarlah \omterm{def:g10:inertia:diagram}{diagram gayanya}.
  \item Bendanya diam, sehingga \omterm{def:g10:inertia:force}{gayanya} \omterm{def:g10:inertia:compensate}{saling meniadakan}
        (\cref{thm:g10:inertia:principle}): bacalah besar yang belum
        diketahui dari gambarnya (panjang yang sama pada satu dimensi,
        komponennya pada dua dimensi).
\end{enumerate}
\end{method}

\begin{example}[Buku di atas meja]\label{ex:g10:inertia:book}
Sebuah buku bermassa \qty{0.50}{kg} terletak di atas meja: \omterm{def:g10:universal-gravitation:weight}{beratnya}
$P = 0.50 \times 9.81 \approx \qty{4.9}{N}$ dan \omterm{def:g10:inertia:inventory}{gaya normal} $\vect R$.
Karena diam, $\vect R$ meniadakan $\vect P$: tegak, ke atas,
$R = \qty{4.9}{N}$. Tambahkan buku kedua dan $R$ menyesuaikan diri ---
sebuah tumpuan mendorong persis sekuat yang diperlukan.
\end{example}

\begin{omfigure}
\begin{tikzpicture}[scale=0.9, font=\small]
  \draw[fill=gray!15] (-2.6,2.2) rectangle (2.6,2.45);
  \draw (0,0) -- (-2.2,2.2);
  \draw (0,0) -- (2.2,2.2);
  \draw[dashed, gray] (0,0) -- (0,2.2);
  \draw[gray] (0,1.0) arc (90:135:1.0);
  \node[gray] at (-0.42,1.38) {$45^\circ$};
  \draw[fill=omExo!25] (0,-0.35) circle (0.35);
  \draw[-Stealth, omDef, very thick] (0,0) -- (-1.06,1.06)
    node[left] {$\vect T_1$};
  \draw[-Stealth, omDef, very thick] (0,0) -- (1.06,1.06)
    node[right] {$\vect T_2$};
  \draw[-Stealth, omThm, very thick] (0,0) -- (0,-1.5)
    node[right] {$\vect P$};
\end{tikzpicture}

{\small Lampu gantung: bagian mendatar kedua tegangannya \omterm{def:g10:inertia:compensate}{saling
meniadakan} karena setangkup; bagian tegaknya bersama-sama memikul
\omterm{def:g10:universal-gravitation:weight}{beratnya}.}
\end{omfigure}

\begin{example}[Lampu pada dua kawat]\label{ex:g10:inertia:lamp}
Sebuah lampu bermassa \qty{3.0}{kg} ($P \approx \qty{29.4}{N}$)
tergantung pada dua kawat, masing-masing bersudut $45^\circ$ dari arah
tegak. Karena setangkup, kedua tegangannya sama ($T_1 = T_2 = T$) dan
komponen mendatarnya \omterm{def:g10:inertia:compensate}{saling meniadakan}; komponen tegaknya,
masing-masing $T \cos 45^\circ$, memikul \omterm{def:g10:universal-gravitation:weight}{beratnya}:
\[
2\,T \cos 45^\circ = P
\qquad\text{sehingga}\qquad
T = \frac{29.4}{2 \times 0.707} \approx \qty{21}{N}
\]
--- lebih dari separuh \omterm{def:g10:universal-gravitation:weight}{beratnya} pada tiap kawat: menarik dengan sudut
memakan tegangan.
\end{example}

\section{Latihan}

\begin{exercise}[$\star$]\label{exo:g10:inertia:1}
Daftarkan dan golongkan (\omterm{def:g10:inertia:contact}{gaya sentuh} atau \omterm{def:g10:inertia:contact}{gaya jarak jauh}) \omterm{def:g10:inertia:force}{gaya} yang
bekerja pada: (a) buku yang diam di atas meja; (b) lampu yang
tergantung pada kawatnya; (c) keping yang meluncur di atas es yang
licin; (d) apel yang jatuh (hambatan udara diabaikan).
\end{exercise}

\begin{exercise}[$\star$]\label{exo:g10:inertia:2}
Hitunglah \omterm{def:g10:universal-gravitation:weight}{berat} sebuah buku \qty{250}{g} dan \omterm{def:g10:universal-gravitation:weight}{berat} seorang manusia
\qty{60}{kg}. Pada skala $\qty{1}{cm} \leftrightarrow \qty{1}{N}$,
berapa panjang anak panah bukunya, dan mengapa skala itu tidak berguna
untuk manusianya?
\end{exercise}

\begin{exercise}[$\star$]\label{exo:g10:inertia:3}
Pada sebuah gambar berskala $\qty{1}{cm} \leftrightarrow \qty{2}{N}$,
sebuah \omterm{def:g10:inertia:force}{gaya} berupa anak panah mendatar sepanjang \qty{3.5}{cm} yang
menunjuk ke kiri, bertitik tangkap di titik $A$. Sebutkan ketiga
cirinya.
\end{exercise}

\begin{exercise}[$\star$]\label{exo:g10:inertia:4}
Sebuah kamus bermassa \qty{1.2}{kg} terletak di atas rak. Gambarlah
\omterm{def:g10:inertia:diagram}{diagram gayanya} dan sebutkan besar serta arah tiap \omterm{def:g10:inertia:force}{gayanya}.
\end{exercise}

\begin{exercise}[$\star$]\label{exo:g10:inertia:5}
Benar atau salah --- perbaikilah yang salah: (a) sebuah benda akhirnya
selalu berhenti kecuali ada \omterm{def:g10:inertia:force}{gaya} yang terus mendorongnya; (b) \omterm{def:g10:inertia:uniform}{gerak
lurus beraturan} tidak memerlukan \omterm{def:g10:inertia:force}{gaya} yang tak terimbangi; (c) benda
yang diam tidak dikenai \omterm{def:g10:inertia:force}{gaya} apa pun; (d) jika kecepatan sebuah benda
berubah, \omterm{def:g10:inertia:force}{gaya} yang bekerja padanya tidak \omterm{def:g10:inertia:compensate}{saling meniadakan}.
\end{exercise}

\begin{exercise}[$\star\star$]\label{exo:g10:inertia:6}
Setelah dilepaskan, sebuah keping menempuh \qty{24}{m} es dalam
\qty{3.0}{s}, lurus dengan kelajuan tetap. Hitunglah lajunya; apakah
\omterm{def:g10:inertia:force}{gaya} yang bekerja padanya \omterm{def:g10:inertia:compensate}{saling meniadakan}? Di atas beton keping yang
sama berhenti dalam beberapa \omterm{def:g10:orders-of-magnitude:unit}{meter}: apa yang berubah pada daftar
\omterm{def:g10:inertia:force}{gayanya}, dan apa yang disimpulkan
\cref{prop:g10:inertia:converse}?
\end{exercise}

\begin{exercise}[$\star\star$]\label{exo:g10:inertia:7}
Seorang penumpang bermassa \qty{70}{kg} berdiri di dalam lift yang naik
dengan laju tetap \qty{1.5}{m/s}. Daftarkan \omterm{def:g10:inertia:force}{gaya} yang bekerja pada
penumpang itu; apakah \omterm{def:g10:inertia:force}{gayanya} \omterm{def:g10:inertia:compensate}{saling meniadakan}? Sebutkan besar \omterm{def:g10:inertia:force}{gaya}
dari lantainya. Apakah ada yang berubah jika liftnya diam?
\end{exercise}

\begin{exercise}[$\star\star$]\label{exo:g10:inertia:8}
Sebuah lampu bermassa \qty{1.8}{kg} tergantung diam pada satu kawat
tegak. Hitunglah tegangan kawat itu. Kawatnya diganti dengan dua kawat
\emph{tegak} yang memikul beban sama rata: berapa tegangan
masing-masing?
\end{exercise}

\begin{exercise}[$\star\star$]\label{exo:g10:inertia:9}
Sebuah bus mengerem mendadak dan penumpang yang berdiri terhuyung ke
depan; pada tikungan dengan kelajuan tetap, penumpang yang sama miring
ke luar. Jelaskan keduanya tanpa mengarang \omterm{def:g10:inertia:force}{gaya} ``ke depan'' atau ``ke
luar''.
\end{exercise}

\begin{exercise}[$\star\star$]\label{exo:g10:inertia:10}
Seorang penerjun payung bermassa seluruhnya \qty{85}{kg} turun tegak
dengan laju tetap \qty{5.0}{m/s}. Gambarlah \omterm{def:g10:inertia:diagram}{diagram gayanya} dan
hitunglah hambatan udara pada payungnya.
\end{exercise}

\begin{exercise}[$\star\star$]\label{exo:g10:inertia:11}
Seorang peselancar air ditarik lurus dengan kelajuan tetap; talinya yang
mendatar menarik dengan \qty{300}{N}. Berapa hambatan mendatar air pada
papan selancarnya (berikan alasannya)? Dua \omterm{def:g10:inertia:force}{gaya} tegak yang mana yang
\omterm{def:g10:inertia:compensate}{saling meniadakan}?
\end{exercise}

\begin{exercise}[$\star\star\star$]\label{exo:g10:inertia:12}
Sebuah lampu bermassa \qty{2.4}{kg} tergantung pada dua kawat bersudut
$40^\circ$ dari arah tegak ($\cos 40^\circ \approx 0.766$). Hitunglah
tiap tegangannya. Mengapa tegangannya membesar ketika kawatnya makin
mendekati arah mendatar, dan mengapa dua kawat yang benar-benar
mendatar tidak akan pernah sanggup menahan lampu itu?
\end{exercise}

\begin{exercise}[$\star\star\star$]\label{exo:g10:inertia:13}
Sebuah batu curling didorong, dilepaskan, meluncur sambil melambat
sangat perlahan, lalu berhenti. Untuk tiap tahapnya (dorongan,
luncuran, diam), gambarlah \omterm{def:g10:inertia:diagram}{diagram gayanya} dan nyatakan apakah \omterm{def:g10:inertia:force}{gayanya}
\omterm{def:g10:inertia:compensate}{saling meniadakan}, dengan mengutip
\cref{thm:g10:inertia:principle} atau \cref{prop:g10:inertia:converse}.
\end{exercise}

\begin{exercise}[$\star\star\star$]\label{exo:g10:inertia:14}
Wahana \emph{Voyager 1} meluncur menembus ruang antarbintang dengan
\qty{17}{km/s}, mesinnya mati. Mengapa ia tidak memerlukan bahan bakar
untuk mempertahankan lajunya, dan apa yang akan diramalkan Aristoteles?
Sejauh apa ia menempuh perjalanan dalam setahun, dalam \unit{km} lalu
dalam \omterm{def:g10:orders-of-magnitude:lightyear}{satuan astronomi} ($\qty{1}{au} \approx
\qty{1.496e8}{km}$)?
Mengapa ia tetap harus menyalakan pendorongnya bahkan sekadar untuk
mengubah arah \emph{tanpa} mengubah lajunya?
\end{exercise}

\begin{exercise}[$\star\star\star$]\label{exo:g10:inertia:15}
Tiga tali menarik sebuah gelang yang diam. Pada sebuah kisi yang satu
kotaknya bernilai \qty{1}{N}, dua \omterm{def:g10:inertia:force}{gayanya} terbaca $\vect F_1\,(3, 2)$
dan $\vect F_2\,(-1, 2)$. Tentukan komponen, besar dan arah \omterm{def:g10:inertia:force}{gaya} yang
ketiga.
\end{exercise}

\section{Soal: Kelembaman dalam tiga babak}

\begin{problem}[{Soal akhir pekan --- jalur berjalan bandara, gelanggang
es dan payung terjun: kecepatan tetap sebagai tanda tangan gaya yang
saling meniadakan, dan mengapa payung terjun tidak mengurangi
hambatannya melainkan lajunya}]
\label{pb:g10:inertia:1}
Tiga pemandangan --- sebuah koper di atas jalur berjalan, sebuah keping
di gelanggang es, seorang penerjun payung di udara senja --- dan satu
hukum yang berbicara pada ketiganya.

\medskip
\noindent\textbf{Bagian I --- Jalur berjalan di bandara.}
Sebuah koper \qty{23}{kg} berdiri di atas jalur berjalan yang bergerak
dengan laju tetap \qty{0.75}{m/s}.
\begin{enumerate}
  \item Bagaimana gerak koper itu di dalam \omterm{def:g10:relative-motion:frames}{kerangka tanah}? Daftarkan
        \omterm{def:g10:inertia:force}{gaya} yang bekerja padanya.
  \item Apakah \omterm{def:g10:inertia:force}{gaya-gaya} itu \omterm{def:g10:inertia:compensate}{saling meniadakan}? Berikan alasannya;
        hitunglah kedua besarnya.
  \item Bandingkan \omterm{def:g10:inertia:diagram}{diagram gayanya} dengan diagram koper yang sama di
        atas lantai. Pernyataan mendalam apa yang diperagakan
        perbandingan itu?
  \item Jalur itu tersentak berhenti. Lukiskan apa yang dilakukan koper
        itu, tanpa mengarang \omterm{def:g10:inertia:force}{gaya} ke depan.
  \item Pemiliknya berjalan \qty{1.2}{m/s} terhadap jalur berjalan itu.
        Berapa lajunya terhadap tanah? Apakah \omterm{def:g10:inertia:force}{gaya} yang bekerja pada
        \emph{dirinya} \omterm{def:g10:inertia:compensate}{saling meniadakan}?
\end{enumerate}

\medskip
\noindent\textbf{Bagian II --- Gelanggang es.}
\begin{enumerate}[resume]
  \item Sebuah keping \qty{160}{g} meluncur lurus dengan laju tetap
        \qty{8.0}{m/s}. Gambarlah \omterm{def:g10:inertia:diagram}{diagram gayanya} beserta besarnya.
  \item Aristoteles mengaku bahwa keping itu memerlukan penggerak. Apa
        jawaban gelanggang esnya, dan mengapa permukaan sehari-hari
        tampak berpihak kepadanya?
  \item Keping itu melintasi bercak yang kasar lalu melambat. Apa
        kesimpulan tentang \omterm{def:g10:inertia:force}{gayanya} di sana? \omterm{def:g10:inertia:force}{Gaya} mana biang keladinya?
  \item Seorang peluncur berbelok mulus dengan kelajuan tetap. Apakah
        \omterm{def:g10:inertia:force}{gaya} yang bekerja padanya \omterm{def:g10:inertia:compensate}{saling meniadakan}? Hati-hati: besaran
        macam apakah kecepatan itu?
  \item Pasangkan diagram dengan geraknya: (a) \omterm{def:g10:inertia:compensate}{gaya saling meniadakan};
        (b) ada \omterm{def:g10:inertia:force}{gaya} tambahan yang berlawanan dengan $\vect v$; (c) ada
        \omterm{def:g10:inertia:force}{gaya} tambahan yang tegak lurus $\vect v$ --- (i) melambat pada
        garis lurus; (ii) berbelok dengan kelajuan tetap; (iii) \omterm{def:g10:inertia:uniform}{gerak
        lurus beraturan}.
\end{enumerate}

\medskip
\noindent\textbf{Bagian III --- Payung terjun.}
Massa seorang penerjun beserta seluruh perlengkapannya adalah
\qty{80}{kg}.
\begin{enumerate}[resume]
  \item Ketika baru meninggalkan pesawat, laju tegak dan hambatannya
        masih nol. Hitunglah \omterm{def:g10:inertia:force}{gayanya}; apakah keduanya \omterm{def:g10:inertia:compensate}{saling
        meniadakan}? Apa yang harus terjadi berikutnya?
  \item Hambatannya tumbuh bersama laju. Lukiskan jatuhnya selama
        hambatannya $<$ \omterm{def:g10:universal-gravitation:weight}{beratnya}.
  \item Sebelum payungnya mengembang, penerjun itu mencapai laju tetap
        \qty{50}{m/s} (``laju terminal''). Hitunglah hambatannya. Gerak
        apa itu, dan mengapa ia disebut ``kelembaman yang menyamar''?
  \item Payungnya mengembang: hambatannya kini jauh melampaui \omterm{def:g10:universal-gravitation:weight}{beratnya}.
        Apa yang terjadi pada kecepatannya? (Penerjunnya bergerak
        \emph{ke bawah} sepanjang waktu.)
  \item Laju tetap yang baru, \qty{5.0}{m/s}, tercapai. Hitunglah
        hambatannya; bandingkan dengan pertanyaan 13.
  \item Puncak ceritanya: apa yang diubah payung terjun itu ---
        hambatan pada keadaan terminalnya, ataukah laju saat
        hambatannya menyamai \omterm{def:g10:universal-gravitation:weight}{beratnya}?
\end{enumerate}

\medskip
\noindent\textbf{Bagian IV --- Putusan.}
\begin{enumerate}[resume]
  \item Lengkapi dan berikan alasannya: ``diam atau bergerak lurus
        beraturan $\iff$ \dots''. Mengapa keadaan diam tidak pantas
        mendapat hukum tersendiri?
  \item Pada babak yang mana ``gerak memerlukan penggerak'' milik
        Aristoteles paling kasatmata gagalnya, dan apa yang akan
        ditunjuk Galileo?
  \item Sebuah lembar ujian menuliskan: ``penerjun payung itu jatuh
        dengan laju tetap, jadi tidak ada \omterm{def:g10:inertia:force}{gaya} yang bekerja padanya''.
        Perbaikilah dalam satu baris.
  \item Penutup: untuk tiap babak, sebutkan pasangan \omterm{def:g10:inertia:force}{gaya} yang \omterm{def:g10:inertia:compensate}{saling
        meniadakan}, lalu nyatakan satu hukum yang dipakai dua puluh
        kali di dalam soal ini.
\end{enumerate}
\end{problem}
